% !TEX root = main.tex
\chapter*{Resum\label{Resum}}
\addcontentsline{toc}{chapter}{Resum}
\fancyhead[RE]{\slshape \textsf{Resum}} %
\fancyhead[LO]{\slshape \textsf{Resum}} %

Els satèl·lits d'observació de la Terra adquireixen volums creixents de dades
multiespectrals sota restriccions d'enllaç de baixada, memòria i processament a
bord. Aquesta memòria implementa en C el pipeline de compressió i descompressió
de SORTENY, un còdec neuronal d'imatges multiespectrals, i l'estén amb control
espacial de qualitat sense selecció manual de regions d'interès.

La qualitat s'assigna per blocs mitjançant mapes generats automàticament a partir
del contingut espectral i espacial de la imatge. Els detectors identifiquen els
blocs prioritaris i les polítiques de qualitat decideixen com es distribueix la
fidelitat entre aquests blocs i el fons. S'han estudiat dues estratègies
complementàries: una prioritza l'estalvi de bits degradant les regions no
prioritàries, i l'altra fixa una qualitat alta dins la regió detectada per
afavorir-ne la conservació.

L'avaluació va incloure 120 retalls Sentinel-2 de vuit bandes. El bitrate
codificat es va mesurar posteriorment amb el pipeline de referència sobre 20
retalls seleccionats. L'estratègia orientada a l'estalvi va reduir el bitrate
mitjà un 8,84\% respecte d'aplicar una qualitat uniforme a tota la imatge,
mentre mantenia o millorava lleugerament la fidelitat de les regions
prioritzades i concentrava la degradació al fons. L'estratègia de preservació va
protegir més aquestes regions a canvi d'un estalvi menor.

La implementació C també es va avaluar en una Raspberry Pi 3B+ com a plataforma
representativa de recursos limitats. Respecte de la versió C inicial executada
amb els mateixos quatre fils, les optimitzacions van reduir el temps total
mitjà un 17,0\%, equivalent a una acceleració d'$1,205\times$, sense canvis
funcionals detectats ni un augment rellevant de memòria. La generació automàtica
del mapa de qualitat va representar menys del 0,05\% del temps del còdec. Per
tant, la decisió espacial té un cost negligible i el principal coll d'ampolla
és l'execució del còdec neuronal. Com a limitacions, el flux binari C encara no
incorpora el codificador final per entropia o aritmètic. La Raspberry s'utilitza
com a referència computacional d'una plataforma a bord de baix recurs; la
validació sobre el hardware definitiu d'una missió queda fora de l'abast.

\textbf{Paraules clau:} compressió d'imatges multiespectrals, compressió
apresa, qualitat fixa, regió d'interès, processament a bord, sistemes encastats,
Raspberry Pi, mapa de qualitat.
